\section{Chapitre~\ref{sec:serocomputer}. Les neurones \nt{SERO}}

Les propriétés des neurones \nt{SERO} eux-mêmes (rythme, autoinhibition, signe fixé par le récepteur, sous-systèmes) sont mesurées directement; les calculs du chapitre (régulateur, filtre, logique) découlent de ses équations; $\tau_c$, le coefficient de diffusion de la sérotonine et le nombre de sites de libération sont des estimations; le rôle de la boucle comme régulateur de niveau dans le cerveau est une hypothèse (dans le noyau du raphé, l'autoinhibition est ponctuelle et ne produit que de courtes pauses), et le rôle de \nt{SERO} comme horizon est une hypothèse appuyée par des expériences comportementales.

{\footnotesize
\setlength{\tabcolsep}{3pt}\renewcommand{\arraystretch}{1.15}
\begin{longtable}{>{\raggedright}p{0.5cm}>{\raggedright}p{4.0cm}>{\raggedright}p{5.6cm}>{\raggedright}p{2.3cm}>{\raggedright\arraybackslash}p{1.7cm}}
\toprule
N\textsuperscript{o} & Affirmation & Expérience et ce qui est mesuré & Source & Statut \\
\midrule\endfirsthead
\toprule
N\textsuperscript{o} & Affirmation & Expérience et ce qui est mesuré & Source & Statut \\
\midrule\endhead
1 & Le signe de l'action de la sérotonine est choisi par le récepteur du destinataire (proposition~\ref{prop:receptor}) & Les récepteurs de la sérotonine se répartissent en familles selon leur pharmacologie et leur mécanisme: 5-HT$_{1A}$ ouvre des canaux potassiques par une protéine G et inhibe la cellule, 5-HT$_{2A}$ et le récepteur ionotrope 5-HT$_3$ l'excitent. Un même neurotransmetteur produit des réponses opposées dans des cellules différentes. & \cite{BarnesSharp1999} & mesuré \\
2 & À l'éveil, le neurone \nt{SERO} émet régulièrement quelques impulsions par seconde & Enregistrement de neurones du noyau du raphé dorsal chez le chat libre de ses mouvements: décharge rythmique lente de $2.8\pm0.2$ imp./s à l'éveil calme; en sommeil lent, la fréquence baisse de $34$--$68\,\%$, en sommeil paradoxal de $98\,\%$. Chez le rat anesthésié, $1.7\pm0.5$ imp./s, très régulièrement. & \cite{TrulsonJacobs1979, GagnonParent2014, JacobsAzmitia1992} & mesuré \\
3 & Le neurone \nt{SERO} est un pacemaker: pas besoin d'une poussée pour chaque impulsion, mais il lui faut un apport de fond régulier (en tranche, la noradrénaline par les récepteurs $\alpha_1$; dans l'organisme, \nt{NORA}) & En tranche de cerveau, les cellules déchargent lentement et régulièrement, chaque impulsion étant suivie d'une post-hyperpolarisation de $200$--$800\ms$. Les cellules spontanément actives sont moins nombreuses en tranche que dans l'organisme: le rythme régulier exige une entrée tonique de noradrénaline par les récepteurs $\alpha_1$, dont le blocage l'éteint. & \cite{VandermaelenAghajanian1983} & mesuré \\
4 & Presque tous les récepteurs sont métabotropes; la réponse est lente: elle monte et retombe en une seconde environ & Toutes les familles de récepteurs, sauf 5-HT$_3$, sont couplées à des protéines G. En tranche, une libération évoquée de sérotonine produit par 5-HT$_{1A}$ un courant inhibiteur qui monte et retombe en moins d'une seconde. & \cite{BarnesSharp1999, CourtneyFord2016} & mesuré \\
5 & Dans les régions cibles, la sérotonine est libérée dans le volume, et pas seulement dans la fente; dans le noyau du raphé lui-même, l'inhibition des voisins se fait point à point & Voltampérométrie cyclique rapide en tranche: la libération par impulsion est la même dans une région riche en synapses et dans le noyau du raphé, où il n'y a presque pas de synapses; elle ne dépend pas du blocage de la recapture ni des récepteurs; il y a bien moins de molécules par terminaison que de transporteurs et de récepteurs, et la concentration extracellulaire coïncide avec l'affinité du récepteur principal. Dans le noyau du raphé, l'inhibition par 5-HT$_{1A}$ se fait point à point: le transporteur empêche la sérotonine de s'accumuler hors de la fente. & \cite{Bunin1998, CourtneyFord2016} & mesuré \\
6 & La concentration selon~\eqref{eq:seroneuron} décroît par recapture; $\tau_c$ de l'ordre de la seconde est une estimation & Voltampérométrie dans le cerveau vivant: la sérotonine extracellulaire est limitée avant tout par la recapture et la dégradation, non par la synthèse. Les travaux cités ne mesurent pas $\tau_c$ directement; l'ordre de grandeur est une estimation du chapitre. & \cite{Hashemi2012serotonin} & mesuré; $\tau_c$: estimation \\
7 & NON et NON-OU sur un seul pacemaker; complétude de la logique \nt{SERO} (proposition~\ref{prop:serocomplete}, fig.~\ref{fig:serogates}) & Découle de~\eqref{eq:seroneuron}: un pacemaker inhibable est un NON-OU, et le NON-OU est fonctionnellement complet. Aucune expérience n'a construit de logique avec des neurones \nt{SERO}; la condition de volumes séparés n'est pas remplie dans le cerveau. & démonstration dans le chapitre & théorie \\
8 & La sérotonine inhibe le neurone \nt{SERO} lui-même par des récepteurs qu'il porte & Chez le rat anesthésié, des agonistes 5-HT$_{1A}$ administrés par voie intraveineuse et par iontophorèse inhibent la décharge des neurones du raphé dorsal avec la même relation dose-réponse que la sérotonine elle-même; les agonistes 5-HT$_{1B}$ n'agissent presque pas. En tranche, la sérotonine endogène des cellules voisines produit par 5-HT$_{1A}$ un courant et une courte pause dans la décharge, sans changer la fréquence moyenne. & \cite{Sprouse1987, CourtneyFord2016} & mesuré \\
9 & Dans le modèle, la boucle par un volume commun est un régulateur de niveau: la dispersion et la constante de temps sont divisées par $1+G$~\eqref{eq:feedback}; que la boucle fonctionne ainsi dans le cerveau est une hypothèse & Formule classique de l'amplificateur à rétroaction négative, redémontrée dans le chapitre. Le gain de boucle $G$ du système \nt{SERO} n'est pas mesuré. Ce qui est mesuré: en tranche, l'autoinhibition se fait point à point, produit une courte pause et ne change pas la fréquence moyenne. & \cite{Black1934feedback, CourtneyFord2016} & théorie; dans le cerveau: hypothèse \\
10 & La concentration est une moyenne glissante de la fréquence, un filtre passe-bas~\eqref{eq:lowpass}, fig.~\ref{fig:lowpass} & Solution de l'équation linéaire~\eqref{eq:seroneuron}; la figure est calculée par cette formule pour $\tau_c=1\,\text{s}$. Pas de modèle séparé dans \texttt{models/}. & démonstration dans le chapitre & théorie \\
11 & La tortuosité des espaces extracellulaires ralentit la diffusion d'une petite molécule d'un facteur $2.6$ environ & Méthode de la source ponctuelle: iontophorèse de tétraméthylammonium et enregistrement de sa concentration par une électrode sélective en tranche et dans le cerveau vivant. Tortuosité $\lambda\approx1.6$, c'est-à-dire un $D$ effectif $\lambda^2\approx2.6$ fois plus petit que le $D$ libre; fraction de volume extracellulaire voisine de $0.2$. Le coefficient de diffusion de la sérotonine elle-même dans le tissu n'est pas mesuré dans ces travaux; dans le chapitre, c'est une estimation. & \cite{Sykova2008} & mesuré \\
12 & La longueur d'un \og fil\fg{} indépendant est $\ell\sim\sqrt{D\tau}\approx20$~\textmu m~\eqref{eq:diffusionlength}; le nombre de fils est limité par le nombre de telles régions (proposition~\ref{prop:volumewires}) & Loi de la diffusion et substitution des estimations de $D$ et $\tau$ (lignes~6 et~11); la proposition~\ref{prop:volumewires} en est une conséquence. Aucune expérience directe ne compte les canaux indépendants de la transmission volumique. & démonstration dans le chapitre & théorie \\
13 & La sortie d'un seul neurone \nt{SERO} couvre d'immenses portions du cerveau; $\sim10^5$ sites de libération selon l'estimation & Reconstruction complète de neurones isolés du raphé dorsal chez le rat: tous les axones ascendants se ramifient abondamment, la longueur totale de l'axone d'une cellule atteint $18.7$~cm, les branches d'une même cellule atteignent le striatum, le cortex préfrontal et l'amygdale. Le nombre de sites de libération par cellule n'a pas été compté directement. & \cite{GagnonParent2014} & mesuré; $10^5$: estimation \\
14 & Les neurones \nt{SERO} se répartissent en quelques sous-systèmes aux sorties et aux réponses différentes & Marquage viro-génétique chez la souris: les neurones qui projettent vers l'amygdale et vers le cortex frontal ne se recouvrent presque pas par leurs ramifications, reçoivent des entrées différentes et répondent de façon opposée à un stimulus désagréable; activer ou inactiver chaque groupe modifie le comportement différemment. & \cite{Ren2018} & mesuré \\
15 & Condition de patience $D<\tau_\gamma\ln(R/r_0)$~\eqref{eq:patience} pour une dépréciation exponentielle; pour la dépréciation mesurée, hyperbolique, $D<\tau_\gamma(R/r_0-1)$ & Découle de la dépréciation $e^{-D/\tau_\gamma}$ ou $1/(1+D/\tau_\gamma)$; démonstration dans le chapitre. Singes, choix avec des délais de quelques secondes: la dépréciation est plus proche de l'hyperbole que de l'exponentielle; la réponse des neurones \nt{DOPA} au signal annonçant une récompense différée décroît elle aussi avec le délai selon une hyperbole. & démonstration dans le chapitre; \cite{KobayashiSchultz2008} & théorie \\
16 & \nt{SERO} fixe l'horizon $\tau_\gamma$ (section~\ref{sec:horizon}) & L'activation optogénétique des neurones \nt{SERO} du raphé dorsal chez la souris allonge l'attente d'une récompense différée et accroît la persévérance dans la recherche d'une récompense devenue rare. Chez l'homme, l'abaissement de la sérotonine (régime sans tryptophane) rend la dépréciation d'une récompense différée plus forte. Comment la concentration se transforme en $\tau_\gamma$, on ne le sait pas. & \cite{Doya2002, Miyazaki2014, Lottem2018, Schweighofer2008} & hypothèse \\
\bottomrule
\end{longtable}}
